\documentclass[12pt]{article}
\usepackage[utf8]{inputenc}
\usepackage[T1]{fontenc}
\usepackage{hyperref}
\title{Palumata Language Resources for K-12 Curriculum Development}
\author{Compiled by Claude Sonnet 5.5 (claude-sonnet-5-5) for Andr\'as Kornai}
\date{}
\begin{document}
\maketitle

\section{Introduction}
Palumata (ISO 639-3 pmc) is described in the sources I found as an extinct and unattested language formerly spoken on Buru Island, Maluku, Indonesia. It is thought to be very closely related to Hukumina, and perhaps a dialect of it. Search results showed no recordings or documentation of the language itself. A curriculum in Palumata is therefore not possible from existing materials.

\section{Primary Educational Resources}
None found.

\section{Linguistic and Reference Materials}
Only an encyclopedic overview was found: \url{https://en.wikipedia.org/wiki/Palumata_language}. It contains no usable lexical or grammatical data.

\section{Religious and Community Resources}
None found.

\section{Digital and Audio Resources}
None found.

\section{Curriculum Development Context}
With no documented lexicon, grammar or texts, no language-medium K-12 curriculum can be built. At most, a heritage-awareness unit on the languages of Buru could be developed, drawing on documented neighbouring languages.
\end{document}